\begin{textsample}{POS Dim 1 – 2015 – Score 23.00 – 2015/t001481.txt}  \label{ex:f1_pos_039}
So I ' ve had \textbf{the} great privilege of traveling to some incredible places, photographing these distant landscapes and remote cultures all over \textbf{the} world.
I love my job.
But people think it ' s this string of epiphanies and sunrises and rainbows, when in reality, it looks more something like this.
This is my office.
We can ' t afford \textbf{the} fanciest places to stay at night, so we tend to sleep a lot outdoors.
As long as we can stay dry, that ' s a bonus.
We also can ' t afford \textbf{the} fanciest restaurants.
So we tend to eat whatever ' s on \textbf{the} local menu.
And if you ' re in \textbf{the} Ecuadorian P{\EmojiFont �}{\EmojiFont �}ramo, you ' re going to eat a large rodent called a cuy.
But what makes our experiences perhaps a little bit different and a little more unique than that of \textbf{the} average person is that we have this gnawing thing in \textbf{the} back of our mind that even in our darkest moments, and those times of despair, we think,"Hey, there might be an image to be made here, there might be a story to be told."And why is storytelling important?
Well, it helps us to connect with our cultural and our natural heritage.
And in \textbf{the} Southeast, there ' s an alarming disconnect between \textbf{the} public and \textbf{the} natural areas that allow us to be here in \textbf{the} first place.
We ' re visual creatures, so we use what we see to teach us what we know.
Now \textbf{the} majority of us aren ' t going to willingly go way down to a swamp.
So how can we still expect those same people to then advocate on behalf of their protection?
We can ' t.
So my job, then, is to use photography as a communication tool, to help bridge \textbf{the} gap between \textbf{the} science and \textbf{the} aesthetics, to get people talking, to get them thinking, and to hopefully, ultimately, get them caring.
I started doing this 15 years ago right here in Gainesville, right here in my backyard.
And I fell in love with adventure and discovery, going to explore all these different places that were just minutes from my front doorstep.
There are a lot of beautiful places to find.
Despite all these years that have passed, I still see \textbf{the} world through \textbf{the} eyes of a child and I try to incorporate that sense of wonderment and that sense of curiosity into my photography as often as I can.
And we ' re pretty lucky because here in \textbf{the} South, we ' re still blessed with a relatively blank canvas that we can fill with \textbf{the} most fanciful adventures and incredible experiences.
It ' s just a matter of how far our imagination will take us.
See, a lot of people look at this and they say,"Oh yeah, wow, that ' s a pretty tree."But I don ' t just see a tree — I look at this and I see opportunity.
I see an entire weekend.
Because when I was a kid, these were \textbf{the} types of images that got me off \textbf{the} sofa and dared me to explore, dared me to go find \textbf{the} woods and put my head \textbf{underwater} and see what we have.
And folks, I ' ve been photographing all over \textbf{the} world and I promise you, what we have here in \textbf{the} South, what we have in \textbf{the} Sunshine State, rivals anything else that I ' ve seen.
But yet our tourism industry is busy promoting all \textbf{the} wrong things.
Before most kids are 12, they ' ll have been to Disney World more times than they ' ve been in a canoe or camping under a starry sky.
And I have nothing against Disney or Mickey ; I used to go there, too.
But they ' re missing out on those fundamental connections that create a real sense of pride and ownership for \textbf{the} place that they call home.
And this is compounded by \textbf{the} issue that \textbf{the} landscapes that define our natural heritage and fuel our aquifer for our drinking water have been deemed as scary and dangerous and spooky.
When our ancestors first came here, they warned,"Stay out of these areas, they ' re haunted.
They ' re full of evil spirits and ghosts."I don ' t know where they came up with that idea.
But it ' s actually led to a very real disconnect, a very real negative mentality that has kept \textbf{the} public disinterested, silent, and ultimately, our environment at risk.
We ' re a state that ' s surrounded and defined by water, and yet for centuries, swamps and wetlands have been regarded as these obstacles to overcome.
And so we ' ve treated them as these second-class ecosystems, because they have very little monetary value and of course, they ' re known to harbor alligators and snakes — which, I ' ll admit, these aren ' t \textbf{the} most cuddly of ambassadors.
So it became assumed, then, that \textbf{the} only good swamp was a drained swamp.
And in fact, draining a swamp to make way for \textbf{agriculture} and development was considered \textbf{the} very essence of conservation not too long ago.
But now we ' re backpedaling, because \textbf{the} more we come to learn about these sodden landscapes, \textbf{the} more secrets we ' re starting to \textbf{unlock} about interspecies relationships and \textbf{the} connectivity of \textbf{habitats}, watersheds and flyways.
Take this bird, for example : this is \textbf{the} prothonotary warbler.
I love this bird because it ' s a swamp bird, through and through, a swamp bird.
They nest and they \textbf{mate} and they breed in these old-growth swamps in these flooded forests.
And so after \textbf{the} \textbf{spring}, after they raise their young, they then fly thousand of miles over \textbf{the} Gulf of Mexico into Central and South America.
And then after \textbf{the} winter, \textbf{the} \textbf{spring} rolls around and they come back.
They fly thousands of miles over \textbf{the} Gulf of Mexico.
And where do they go?
Where do they land?
Right back in \textbf{the} same tree.
That ' s \textbf{nuts}.
This is a bird \textbf{the} size of a tennis ball — I mean, that ' s crazy!
I used a GPS to get here today, and this is my hometown.
It ' s crazy.
So what happens, then, when this bird flies over \textbf{the} Gulf of Mexico into Central America for \textbf{the} winter and then \textbf{the} \textbf{spring} rolls around and it flies back, and it comes back to this : a freshly sodded \textbf{golf} course?
This is a narrative that ' s all too commonly unraveling here in this state.
And this is a natural process that ' s \textbf{occurred} for thousands of years and we ' re just now learning about it.
So you can imagine all else we have to learn about these landscapes if we just preserve them first.
Now despite all this rich life that abounds in these swamps, they still have a bad name.
Many people feel uncomfortable with \textbf{the} idea of wading into Florida ' s blackwater.
I can understand that.
But what I loved about growing up in \textbf{the} Sunshine State is that for so many of us, we live with this latent but very palpable fear that when we put our toes into \textbf{the} water, there might be something much more ancient and much more adapted than we are.
Knowing that you ' re not top dog is a welcomed discomfort, I think.
How often in this modern and urban and digital age do you actually get \textbf{the} chance to feel vulnerable, or consider that \textbf{the} world may not have been made for just us?
So for \textbf{the} last decade, I began seeking out these areas where \textbf{the} concrete yields to forest and \textbf{the} pines turn to cypress, and I viewed all these mosquitoes and reptiles, all these discomforts, as affirmations that I ' d found true wilderness, and I embrace them wholly.
Now as a conservation photographer obsessed with blackwater, it ' s only fitting that I ' d eventually end up in \textbf{the} most \textbf{famous} swamp of all : \textbf{the} Everglades.
Growing up here in North Central Florida, it always had these enchanted names, places like Loxahatchee and Fakahatchee, Corkscrew, Big Cypress.
I started what turned into a five-year project to hopefully reintroduce \textbf{the} Everglades in a new light, in a more inspired light.
But I knew this would be a tall order, because here you have an area that ' s roughly a third \textbf{the} size \textbf{the} state of Florida, it ' s huge.
And when I say Everglades, most people are like,"Oh, yeah, \textbf{the} national park."But \textbf{the} Everglades is not just a park ; it ' s an entire watershed, starting with \textbf{the} Kissimmee chain of lakes in \textbf{the} north, and then as \textbf{the} rains would fall in \textbf{the} summer, these downpours would flow into Lake Okeechobee, and Lake Okeechobee would fill up and it would overflow its banks and \textbf{spill} southward, ever slowly, with \textbf{the} topography, and get into \textbf{the} river of grass, \textbf{the} Sawgrass Prairies, before meting into \textbf{the} cypress slews, until going further south into \textbf{the} mangrove swamps, and then finally — finally — reaching Florida Bay, \textbf{the} emerald gem of \textbf{the} Everglades, \textbf{the} great estuary, \textbf{the} 850 square-mile estuary.
So sure, \textbf{the} national park is \textbf{the} southern end of this system, but all \textbf{the} things that make it unique are these inputs that come in, \textbf{the} fresh water that starts 100 miles north.
So no manner of these political or invisible boundaries protect \textbf{the} park from polluted water or insufficient water.
And unfortunately, that ' s precisely what we ' ve done.
Over \textbf{the} last 60 years, we have drained, we have dammed, we have dredged \textbf{the} Everglades to where now only one third of \textbf{the} water that used to reach \textbf{the} \textbf{bay} now reaches \textbf{the} \textbf{bay} today.
So this story is not all sunshine and rainbows, unfortunately.
For better or for worse, \textbf{the} story of \textbf{the} Everglades is intrinsically tied to \textbf{the} peaks and \textbf{the} valleys of mankind ' s relationship with \textbf{the} natural world.
But I ' ll show you these beautiful pictures, because it gets you on board.
And while I have your attention, I can tell you \textbf{the} real story.
It ' s that we ' re taking this, and we ' re trading it for this, at an alarming rate.
And what ' s lost on so many people is \textbf{the} \textbf{sheer} scale of which we ' re discussing.
Because \textbf{the} Everglades is not just responsible for \textbf{the} drinking water for 7 million Floridians ; today it also provides \textbf{the} agricultural fields for \textbf{the} year-round tomatoes and oranges for over 300 million Americans.
And it ' s that same seasonal \textbf{pulse} of water in \textbf{the} summer that built \textbf{the} river of grass 6, 000 years ago.
Ironically, today, it ' s also responsible for \textbf{the} over half a million acres of \textbf{the} endless river of sugarcane.
These are \textbf{the} same fields that are responsible for dumping exceedingly high levels of fertilizers into \textbf{the} watershed, forever changing \textbf{the} system.
But in order for you to not just understand how this system works, but to also get personally connected to it, I decided to break \textbf{the} story down into several different narratives.
And I wanted that story to start in Lake Okeechobee, \textbf{the} beating heart of \textbf{the} Everglade system.
And to do that, I picked an ambassador, an iconic species.
This is \textbf{the} Everglade snail kite.
It ' s a great bird, and they used to nest in \textbf{the} thousands, thousands in \textbf{the} northern Everglades.
And then they ' ve gone down to about 400 nesting pairs today.
And why is that?
Well, it ' s because they eat one source of food, an apple snail, about \textbf{the} size of a ping-pong ball, an aquatic gastropod.
So as we started damming up \textbf{the} Everglades, as we started diking Lake Okeechobee and draining \textbf{the} wetlands, we lost \textbf{the} \textbf{habitat} for \textbf{the} snail.
And thus, \textbf{the} population of \textbf{the} kites declined.
And so, I wanted a photo that would not only communicate this relationship between wetland, snail and bird, but I also wanted a photo that would communicate how incredible this relationship was, and how very important it is that they ' ve come to depend on each other, this healthy wetland and this bird.
And to do that, I brainstormed this idea.
I started sketching out these plans to make a photo, and I sent it to \textbf{the} wildlife biologist down in Okeechobee — this is an endangered bird, so it takes special permission to do.
So I built this submerged platform that would hold snails just right under \textbf{the} water.
And I spent months planning this crazy idea.
And I took this platform down to Lake Okeechobee and I spent over a week in \textbf{the} water, wading waist-deep, 9-hour shifts from dawn until dusk, to get one image that I thought might communicate this.
And here ' s \textbf{the} day that it finally worked : [ Video : After setting up \textbf{the} platform, I look off and I see a kite coming over \textbf{the} cattails.
And I see him scanning and \textbf{searching}.
And he gets right over \textbf{the} trap, and I see that he ' s seen it.
And he beelines, he goes straight for \textbf{the} trap.
And in that moment, all those months of planning, waiting, all \textbf{the} sunburn, mosquito \textbf{bites} — suddenly, they ' re all worth it.
Oh my gosh, I can ' t believe it! ] You can believe how excited I was when that happened.
But what \textbf{the} idea was, is that for someone who ' s never seen this bird and has no reason to care about it, these photos, these new perspectives, will help shed a little new light on just one species that makes this watershed so incredible, so valuable, so important.
Now, I know I can ' t come here to Gainesville and talk to you about animals in \textbf{the} Everglades without talking about gators.
I love gators, I grew up loving gators.
My parents always said I had an unhealthy relationship with gators.
But what I like about them is, they ' re like \textbf{the} freshwater equivalent of sharks.
They ' re feared, they ' re hated, and they are tragically misunderstood.
Because these are a unique species, they ' re not just apex predators.
In \textbf{the} Everglades, they are \textbf{the} very architects of \textbf{the} Everglades, because as \textbf{the} water drops down in \textbf{the} winter during \textbf{the} dry season, they start excavating these holes called gator holes.
And they do this because as \textbf{the} water drops down, they ' ll be able to stay \textbf{wet} and they ' ll be able to forage.
And now this isn ' t just affecting them, other animals also depend on this relationship, so they become a keystone species as well.
So how do you make an apex predator, an ancient reptile, at once look like it dominates \textbf{the} system, but at \textbf{the} same time, look vulnerable?
Well, you wade into a pit of about 120 of them, then you hope that you ' ve made \textbf{the} right decision.
I still have all my fingers, it ' s cool.
But I understand, I know I ' m not going to rally you guys, I ' m not going to rally \textbf{the} troops to"Save \textbf{the} Everglades for \textbf{the} gators!"It won ' t happen because they ' re so ubiquitous, we see them now, they ' re one of \textbf{the} great conservation success stories of \textbf{the} US.
But there is one species in \textbf{the} Everglades that no matter who you are, you can ' t help but love, too, and that ' s \textbf{the} roseate spoonbill.
These birds are great, but they ' ve had a really tough time in \textbf{the} Everglades, because they started out with thousands of nesting pairs in Florida Bay, and at \textbf{the} turn of \textbf{the} 20th century, they got down to two — two nesting pairs.
And why?
That ' s because women thought they looked better on their hats then they did flying in \textbf{the} sky.
Then we banned \textbf{the} plume trade, and their numbers started rebounding.
And as their numbers started rebounding, scientists began to pay attention, they started studying these birds.
And what they found out is that these birds ' behavior is intrinsically tied to \textbf{the} annual draw-down cycle of water in \textbf{the} Everglades, \textbf{the} thing that defines \textbf{the} Everglades watershed.
What they found out is that these birds started nesting in \textbf{the} winter as \textbf{the} water drew down, because they ' re tactile feeders, so they have to touch whatever they eat.
And so they wait for these concentrated pools of \textbf{fish} to be able to feed enough to feed their young.
So these birds became \textbf{the} very icon of \textbf{the} Everglades — an indicator species of \textbf{the} overall health of \textbf{the} system.
And just as their numbers were rebounding in \textbf{the} mid-20th century — shooting up to 900, 1, 000, 1, 100, 1, 200 — just as that started happening, we started draining \textbf{the} southern Everglades.
And we stopped \textbf{two-thirds} of that water from moving south.
And it had drastic consequences.
And just as those numbers started reaching their peak, unfortunately, today, \textbf{the} real spoonbill story, \textbf{the} real photo of what it looks like is more something like this.
And we ' re down to less than 70 nesting pairs in Florida Bay today, because we ' ve disrupted \textbf{the} system so much.
So all these different organizations are shouting, they ' re screaming,"\textbf{The} Everglades is fragile!
It ' s fragile!"It is not.
It is resilient.
Because despite all we ' ve taken, despite all we ' ve done and we ' ve drained and we ' ve dammed and we ' ve dredged it, pieces of it are still here, waiting to be put back together.
And this is what I ' ve loved about South Florida, that in one place, you have this unstoppable force of mankind meeting \textbf{the} immovable \textbf{object} of tropical nature.
And it ' s at this new frontier that we are forced with a new appraisal.
What is wilderness worth?
What is \textbf{the} value of biodiversity, or our drinking water?
And fortunately, after decades of debate, we ' re finally starting to act on those questions.
We ' re slowly \textbf{undertaking} these projects to bring more freshwater back to \textbf{the} \textbf{bay}.
But it ' s up to us as citizens, as residents, as stewards to hold our elected officials to their promises.
What can you do to help?
It ' s so easy.
Just get outside, get out there.
Take your friends out, take your kids out, take your family out.
Hire a \textbf{fishing} guide.
Show \textbf{the} state that protecting wilderness not only makes ecological sense, but economic sense as well.
It ' s a lot of fun, just do it — put your feet in \textbf{the} water. \textbf{The} swamp will change you, I promise.
Over \textbf{the} years, we ' ve been so generous with these other landscapes around \textbf{the} country, cloaking them with this American pride, places that we now consider to define us : Grand Canyon, Yosemite, Yellowstone.
And we use these parks and these natural areas as beacons and as cultural compasses.
And sadly, \textbf{the} Everglades is very commonly left out of that conversation.
But I believe it ' s every bit as iconic and emblematic of who we are as a country as any of these other wildernesses.
It ' s just a different kind of wild.
But I ' m encouraged, because maybe we ' re finally starting to come around, because what was once deemed this swampy wasteland, today is a World Heritage site.
It ' s a wetland of international importance.
And we ' ve come a long way in \textbf{the} last 60 years.
And as \textbf{the} world ' s largest and most ambitious wetland restoration project, \textbf{the} international spotlight is on us in \textbf{the} Sunshine State.
Because if we can heal this system, it ' s going to become an icon for wetland restoration all over \textbf{the} world.
But it ' s up to us to decide which legacy we want to attach our flag to.
They say that \textbf{the} Everglades is our greatest test.
If we pass it, we get to keep \textbf{the} \textbf{planet}.
I love that quote, because it ' s a challenge, it ' s a prod.
Can we do it?
Will we do it?
We have to, we must.
But \textbf{the} Everglades is not just a test.
It ' s also a gift, and ultimately, our responsibility.
Thank you.

% matched lemmas: agriculture, bay, bite, famous, fish, fishing, golf, habitat, mate, nut, object, occur, planet, pulse, search, sheer, spill, spring, the, two-third, undertake, underwater, unlock, wet
\end{textsample}
